\section*{Executive summary}

The UK public sector is rich in data (open data through to highly sensitive operational and personal data), but it is often \emph{hard to turn into answers}. \cite{oecd_data_driven_public_sector2019,nao_data_across_gov} Historically, we built applications as the primary interface to data, and trained people to use the applications.

This brief proposes a complementary model: \textbf{Apps $\rightarrow$ Answers}. In this model, people ask questions in plain language; an intelligence layer (increasingly AI) interprets intent; and a controlled ``middle layer'' transforms raw data into usable information with provenance, permissions, and auditability. \cite{mcp_spec_2025_11_25,prov_o,prov_dm}

The middle layer is where practical responsibility sits: (1) \textbf{dataset scoping and permissions}; (2) \textbf{discovery and semantics} (what data exists, what it means); (3) \textbf{provenance and audit} (how an answer was produced, and how to reproduce it); (4) \textbf{safety guardrails} (policy, privacy, and acceptable use); and (5) \textbf{operations} (cost, latency, reliability, and degradation). \cite{ncsc_secure_ai_guidelines,ico_ai_data_protection,data_ai_ethics_framework,five_safes_framework,algorithmic_transparency_recording_standard}

We use \textbf{Model Context Protocol (MCP)} and an example server (\textbf{MCP-Geo}, combining Ordnance Survey and Office for National Statistics data) to make this concrete: \emph{capabilities over applications}. The case study is also used to show how errors and ``unknowns'' should surface (with provenance) and how audit traces support assurance. \cite{mcp_spec_2025_11_25,prov_o,algorithmic_transparency_recording_standard,artefact_claude_stopped_after_ngd_features_2026_02_18,artefact_codex_trace_jsonl_2026_02_18}

We focus on selecting initial use cases that are repeatable, bounded in scope, and measurable, then adopting safely in stages: \textbf{open} (public datasets) $\rightarrow$ \textbf{controlled} (internal data with role-based access) $\rightarrow$ \textbf{sensitive} (high-impact contexts with tighter policy, monitoring, and human oversight). \cite{five_safes_framework,ico_ai_data_protection,ncsc_secure_ai_guidelines}

Finally, we propose a 30-day pilot ``thin slice'': one priority question-set, one constrained dataset bundle, one middle-layer implementation (MCP-backed), and one auditable end-to-end workflow that proves value while making risk visible.

\begin{itemize}[leftmargin=*]
  \item \textbf{What you’ll leave with:} a clear explanation of Apps $\rightarrow$ Answers and why it matters.
  \item A concrete checklist of middle-layer responsibilities (permissions, semantics, provenance, guardrails, operations).
  \item A simple lens for choosing strong first use cases (repeatable, bounded risk, measurable value).
  \item A staged adoption path (open $\rightarrow$ controlled $\rightarrow$ sensitive) with explicit guardrails.
  \item A 30-day pilot plan for an auditable ``thin slice'' you can run with a small team.
\end{itemize}
